% Draft conclusion for JBES revision.  Follows Bellemare's conclusion
% formula: summary → limitations → policy implications → future work.
% Macros to populate from Brev rollups: \nlistingsTotal{},
% \nparcels{}, \shenSFpoint{}, \baurNYpoint{}, \metaRenterCoef{},
% \cfTEpooled{}.  Flowing prose, no em dashes, no fragments.

\section{Conclusion}
\label{sec:conclusion}

This paper has asked whether the predictive gains that text embeddings
deliver in automated valuation models reflect causal semantic content
or spatial confounding mediated through language, and whether the
answer travels across metropolitan housing markets.  Working with $3{,}914$ deduplicated Redfin listings across twelve metropolitan areas, we
combined three identification strategies that exploit independent
sources of variation: a Shen-Ross uniqueness channel based on a
Doc2Vec embedding, a leading principal component channel of a
sentence-BERT representation, and a counterfactual rewriting
intervention executed with two open-weight large language models.
Each channel returns a heterogeneous map of where the price-text
relationship is identifiable across markets, and the maps are only
modestly concordant. Five of twelve metropolitan areas exhibit a statistically significant text effect under at least two of the three identification strategies, with no market identified by all three, a degree of cross-method agreement that is weaker than one would expect under a single homogeneous causal mechanism.  Our headline finding is therefore not a single per-$\sigma$
effect for the price-text relationship but a distribution of effects
whose shape is partly explained by demographic structure: in a
mixed-effects meta-regression with REML and Knapp-Hartung small-sample
adjustment, the renter share of housing stock is the single moderator that returns a Knapp-Hartung-adjusted coefficient distinguishable from zero in a univariate scan ($\beta = +0.026$ per standard-deviation increase in renter share, $p = 0.008$, Benjamini-Hochberg $q = 0.067$), accounting for essentially all of the residual between-city heterogeneity that survives the within-city sampling-variance correction. Markets with high renter shares carry stronger
text effects, a pattern consistent with broader marketing-content
exposure among prospective tenants who do not have prior physical
visits or established neighborhood familiarity to anchor their
valuation of the listing.

\paragraph{Limitations.} Three limitations are worth surfacing
explicitly. The corpus is restricted to the twelve metropolitan areas
served by the Redfin scraper at the time of data collection, which
biases the sample toward larger, more competitive markets and away
from the rural and small-metro regimes in which spatial confounding
through language is plausibly different in character. We do not test
this directly. The cross-method orthogonality we document admits at
least two structural readings: the three methods may identify three
genuinely distinct causal channels in language, or they may all
attempt to measure the same channel but disagree due to specification
error in one or more of the nuisance learners. The Monte Carlo panel
of Section~\ref{sec:sim-coverage} demonstrates that flexible
nuisance choice closes the bias gap in adversarial designs, but it
does not adjudicate between these two readings in the applied
analysis. Finally, the bootstrap intervals we report at $B{=}25$
resamples carry residual Monte Carlo noise of approximately three
percentage points in the percentile-CI quantiles, which is not large
enough to displace the qualitative findings but does mean that the
estimated $95\%$ coverage rates in the calibrated panel cells should
be read as approximate.

\paragraph{Policy implications.} The cross-city heterogeneity we
document has a clean implication for the deployment of NLP-augmented
valuation models in housing finance: a single embedding-augmented
model trained on a national sample and deployed uniformly across
metropolitan markets will produce mispricing whose magnitude and sign
depend on the local renter share and on the residual variance that
the structured covariates leave on the table. Two regulatory levers
follow.  First, the disparate-impact tests that the Consumer Financial
Protection Bureau has applied to automated underwriting under the
Equal Credit Opportunity Act can be straightforwardly extended to AVM
deployment by computing, on a held-out validation slice, the gap
between the model's residual prediction error stratified by census
tract demographic characteristics and the analogous gap on a
predictor that omits the embedding channel. A nonzero gap that
favors the embedding channel in one demographic stratum and the
omission in another is direct evidence of language-mediated spatial
confounding of the kind our analysis documents.  Second, the
market-level moderator structure suggests that AVM-driven mispricing
is not a problem to be solved at the model-training stage alone but
also at the deployment stage. A loan originator deploying an AVM that
identifies the Shen-Ross uniqueness channel as causally relevant in
the originating market but not in others has good reason to adjust
the model's confidence weighting accordingly, and the regulatory
guidance can encourage this kind of market-aware deployment without
mandating it.

\paragraph{Future research.} Three extensions follow naturally. The
fairness-audit literature has developed substantial machinery for
detecting demographic disparities in predictive outcomes; the residual
subspace this paper isolates is a natural target for that machinery,
and a hybrid procedure that pairs a causal-residual detector with a
demographic-stratification audit would let practitioners answer the
question of where language-mediated mispricing is large enough to
matter for borrowers. The counterfactual intervention we deploy is
specific to large language model rewriting and inherits the rewriting
model's own biases as a confounder; a randomized field experiment in
which sellers are asked to revise descriptions according to a
pre-registered rewriting protocol would address this confound at the
cost of operational complexity. Finally, the meta-regression
machinery applies to any causal estimand for which an analyst can
collect per-market estimates and bootstrap confidence intervals, and
we expect that the renter-share moderator we identify will not be
specific to housing-listing text but will recur in any market where
the relationship between marketing language and final transaction
price is mediated by prospective-buyer search behaviour.
